\bookchapter{第二十九章：デジタルヘルス——スマホとの健康的な関係づくり}

\chapterintro{こんな経験はないか：5分だけ短動画を見るつもりが、気がつくと2時間経っていた？夜ベッドでスマホを見ていたら、見れば見るほど目が冴えて、結局深夜2時に寝た？スマホは私たちの生活を変えたが、静かに私たちの時間と健康をむさぼり食っている。この章はスマホを“やめる”ことが目的ではなく、スマホに対する主導権を取り戻すことが目的だ。}

\booksection{一、まず“スクリーンタイム”を見よう——数字が目を覚まさせる}

変える前に、まず現状を把握しよう。スマホには多くの人が一度も開いたことのない機能がある。“スクリーンタイム”（iPhone）または“デジタルウェルビーイング”（Android）と呼ばれ、毎日スマホにどれだけ時間を使ったか、その時間がどのAppに使われたかを教えてくれる。

\begin{itemize}
  \item \textbf{iPhone}：設定 > スクリーンタイム > すべてのアクティビティを表示。毎日のロック解除回数、スマホを手に取った回数、各Appの使用時間が分かる。
  \item \textbf{Android}：設定 > デジタルウェルビーイングとペアレンタルコントロール。Huaweiは"健康的なスマホ使用"、Xiaomiは"スクリーンタイム管理"と呼ぶ。機能は類似し、毎日の使用データが見られる。
\end{itemize}

\textbf{初めてこのデータを見る人の80\%は驚く}。ある友人はデータを見た後、毎日TikTokに4時間以上使っていることに気づき、一週間で28時間——アルバイト一台分に相当する。この“数字の衝撃”が変化の第一歩だ。

\booksection{二、Appに“使用制限”を設定する——自分にスピードリミッターを付ける}

問題が分かったら、ターゲットを絞って制限を設定できる。Appを一括削除する必要はなく、“一日の枠”を設定するだけでいい：

\begin{enumerate}
  \item \textbf{iPhone}：設定 > スクリーンタイム > App制限 > 制限を追加。制限したいApp（例：TikTok）を選び、毎日30分に設定。時間が来るとAppアイコンがグレーアウトし、開くには追加でパスワード入力が必要。
  \item \textbf{Android}：設定 > デジタルウェルビーイング > Appタイマー。Appを選んで毎日の使用上限を設定、時間が来るとAppアイコンがグレーアウト。
  \item \textbf{テクニック}：制限のパスワードを普段使わない覚えにくいものにする。簡単に入力できると毎回"延長"してしまい、制限が形骸化する。
\end{enumerate}

\booksection{三、“おやすみモード”の正しい使い方——通知を消すのではなく不安を消す}

多くの人のスマホは一日に何十回も通知で鳴り、鳴るたびに注意力が中断される。“ちょっと通知を見るだけ”だと思っていても、実際は中断後に注意力を取り戻すのに15-20分かかる。おやすみモードは“連絡絶ち”ではなく、注意力を取り戻すためのものだ。

\begin{itemize}
  \item \textbf{時間指定でオン}：おやすみモードを夜10時から朝7時まで自動オンに設定。これで睡眠中に迷惑SMSで起こされることがなくなる。
  \item \textbf{重要な連絡先を通過させる}：おやすみモードの設定で、"着信を許可"を"お気に入りの連絡先"または"繰り返し着信"（3分以内の2回目の着信は通常通り鳴る）に設定。これで家族の緊急電話を見逃さず、営業電話にも邪魔されない。
  \item \textbf{仕事集中モード}：iPhoneの"集中モード"はカスタムシーンを設定できる——例えば"仕事"モードでは仕事関連Appの通知のみ許可、"読書"モードではすべての通知をミュート。
\end{itemize}

\booksection{四、目の保護はAppではなく習慣——20-20-20ルール}

長時間スマホ画面を見ると、目の乾きや疲れは避けられない。アイケアモード（ブルーライトカット）はある程度役立つが、本当に効果的なのは目の使い方を変えることだ：

\begin{tcbinfo}{20-20-20ルール}
画面を見るごとに\textbf{20分}、視線を\textbf{20フィート（約6メートル）}以遠に移し、\textbf{20秒}維持する。この簡単な動作が毛様体筋の疲労を効果的に和らげる。

\end{tcbinfo}

\begin{itemize}
  \item \textbf{アイケアモードをオンにする}：設定 > 表示 > アイケアモード（またはNight Shift）、オンにすると画面の色温度が暖色になり、ブルーライトが減少。日没から日の出まで自動オンに設定することを推奨。
  \item \textbf{画面の明るさを調整する}：周囲が暗い時は画面も暗くする。多くの人が暗闇でスマホを見る時明るさを上げすぎるが、これが最も目を痛める行為だ。"自動明るさ"をオンにし、正常に動作していることを確認。
  \item \textbf{距離を保つ}：スマホと目の距離は最低30センチ。多くの人がスマホをとても近くに持って見るが、こうすると毛様体筋が常に緊張収縮状態になり、近視の進行が早い。
\end{itemize}

\booksection{五、“スマホ首”の頸椎自救ガイド}

下を向いてスマホを見る時、首が受ける圧力は通常立っている時の5-6倍だ。長期間下を向くと頸椎のカーブがなくなり、骨棘形成、めまい、手のしびれを引き起こす。以下は簡単な改善方法：

\begin{itemize}
  \item \textbf{スマホを高く上げる}：スマホを見る時はできるだけ目の高さまで上げる。膝の上に置いたスマホを下を向いて見ない。腕は疲れるが、首は感謝する。
  \item \textbf{30分ごとに首を動かす}：ゆっくり左に首を回して限界まで、3秒キープ、正面に戻す；次に右に回し、3秒キープ。5回繰り返す。次に"米字操"——顎で空中に"米"の字を書く。
  \item \textbf{壁に寄りかかって3分立つ}：後頭部、肩甲骨、臀部、かかとの四点を壁に付け、毎日3分立つ。猫背·下向きの姿勢を効果的に改善できる。
\end{itemize}

\booksection{六、睡眠とスマホ——寝室で最も変えるべき習慣}

スマホ画面のブルーライトはメラトニン分泌を抑制し、“見れば見るほど目が冴える”状態にする。多くの人は“見ているうちに眠くなった”と思うが、実際はブルーライトで脳が騙されており、スマホを置いた後かえって眠りにくい。以下は改善のアドバイス：

\begin{itemize}
  \item \textbf{寝る1時間前はスマホを見ない}：最も効果的だが最も難しい一つ。できない場合は、少なくとも"Night Shift/アイケアモード"をオンにしてブルーライトを最低限に。
  \item \textbf{スマホを手の届かない場所に置く}：充電器を寝室の反対側に置く。ベッドサイドではなく。これで思わず手に取って見てしまうのを防ぎ、アラームが鳴っても起きて止めなければならない——一石二鳥。
  \item \textbf{"睡眠モード"をオンにする}：iPhoneの"睡眠"集中モードは睡眠時間帯に自動で画面をロックし、ロック画面をシンプルにする；Android各ブランドにも類似機能あり。
  \item \textbf{スマホの代わりに目覚まし時計を使う}：十数元の物理的目覚まし時計を買い、スマホをアラームにしない。"時間を見るついでに思わずスクロール"の罠がなくなることに気づく。
\end{itemize}

\booksection{七、情報不安症——“選択的盲目”を身につける}

私たちが毎日受け取る情報量は、昔の人が一生で受け取る量だ。しかし人間の脳はこれほど多くの情報を処理するようにはできておらず、情報過多は不安、注意力散漫、意思決定疲労を引き起こす。解決策は“より多くの情報を得る”ことではなく“能動的に遮断する”ことだ：

\begin{itemize}
  \item \textbf{フォローを解除する}：フォローしているすべての公式アカウント、UP主、ブロガーのリストを開き、自問する："過去一ヶ月で実用的な価値を得たか？"なければフォロー解除。情報チャネルは多さより質。
  \item \textbf{不要なAppの通知をオフにする}：WeChat、電話、SMS、カレンダー以外のApp通知をすべてオフに。ニュースは通知不要——自分から見に行けば自然に目に入る。
  \item \textbf{"情報断食日"を設ける}：毎週一日（例：日曜）を選び、終日スマホでニュースを見ず、短動画をスクロールしない。空が崩れ落ちないこと、そして月曜の精神的状態が明らかに良くなることに気づく。
\end{itemize}

\begin{tcbpractice}{実践：一週間デジタルヘルスチャレンジ}
\begin{enumerate}
  \item 今日：スクリーンタイム/デジタルウェルビーイングを開き、先週の使用データを見る。
  \item 明日：最も時間を費やしたエンタメAppに毎日の制限を設定（30-60分を推奨）。
  \item 明後日：おやすみモードを時間指定オン（22:00-07:00）に設定、お気に入り連絡先の着信を許可。
  \item 四日目：寝る1時間前はスマホを見ない、充電器を寝室の反対側に置く。
  \item 五日目：もう見ていない5つの公式アカウントやブロガーのフォローを解除。
  \item 六日目：丸一日"情報断食"に挑戦——ニュースを見ず短動画をスクロールしない。
  \item 七日目：今週のスクリーンタイムレポートを見て、先週と比較する。
\end{enumerate}

\end{tcbpractice}

\begin{tcbsummary}{本章のまとめ}
\begin{itemize}
  \item まずスクリーンタイムデータを見て現状を認識することが変化の第一歩
  \item エンタメAppに制限を設定し、おやすみモードを時間指定でオンにして注意力を取り戻す
  \item 20-20-20ルールは最もシンプルで効果的なアイケア習慣
  \item 下を向いてスマホを見る時の首への負担は5-6倍、スマホを高く上げる+定期的に首を動かす
  \item 寝る1時間前はスマホを見ない、充電器を手の届かない場所に置く
  \item 情報不安の特効薬は"能動的に遮断"することであり、"より多く得る"ことではない
\end{itemize}

\end{tcbsummary}
